(1) 欧拉方程为 $I_1\dot\omega_1=(I_2-I_3)\omega_2\omega_3$、$I_2\dot\omega_2=(I_3-I_1)\omega_3\omega_1$、$I_3\dot\omega_3=(I_1-I_2)\omega_1\omega_2$。无力矩时动能 $E=\frac12\sum_iI_i\omega_i^2$ 与角动量长度平方 $L^2=\sum_i I_i^2\omega_i^2$ 都守恒。
(2) 在中间轴解附近令 $\omega_1=\delta_1,\omega_2=\Omega+O(\delta^2),\omega_3=\delta_3$，保留一阶项，有 $I_1\dot\delta_1=(I_2-I_3)\Omega\delta_3$、$I_3\dot\delta_3=(I_1-I_2)\Omega\delta_1$；$\dot\omega_2$ 只从扰动乘积开始，属于二阶小量。
(3) 消去 $\delta_3$ 得 $\ddot\delta_1=\Omega^2[(I_3-I_2)(I_2-I_1)/(I_1I_3)]\,\delta_1$。因惯量严格递增，括号为正，故增长率 $\gamma=|\Omega|\sqrt{(I_3-I_2)(I_2-I_1)/(I_1I_3)}$，一般扰动含 $e^{\gamma t}$，表明中间轴匀转在微扰下不稳定。
(4) 绕 $I_1$ 轴以 $\Omega$ 转动时其余扰动满足谐振方程，$\omega_{\rm stab,1}=|\Omega|\sqrt{(I_2-I_1)(I_3-I_1)/(I_2I_3)}$；绕 $I_3$ 轴时 $\omega_{\rm stab,3}=|\Omega|\sqrt{(I_3-I_1)(I_3-I_2)/(I_1I_2)}$，均为实频率。能量和角动量守恒将运动限制在其交线，但交线在中间轴附近呈双曲鞍形，仍容许偏离该方向。
